\documentclass[11pt]{article}
\usepackage{mathblatt}
% Präambel des Setzers (werkzeuge/setzer-praeambel.tex, wortgleich); dazu \gteil (Muster T2U)
\makeatletter
\setlength{\mbpfbe}{0mm}% keine Punkte (pruefung.md § 5: nur beim Prüfstein)
\newlength{\gEin}\setlength{\gEin}{0pt}
\let\gSchrift\relax
\renewcommand{\pfkreuzzeile}[1]{\par\noindent\hangindent3.2em\hangafter1\hspace*{1.6em}$\square$\ #1\par}
\newcommand{\gkreuz}[1]{$\square$\ #1\hspace{2em}}
\newcommand{\gkreuzzeile}[2]{\par\noindent\hangindent2.2em\hangafter1\mbox{#1\,$\square$}\ #2\par}% Buchstabe, Option: eine je Zeile (neben Skizze)
% links, was man ansieht, rechts, was man tut; Spaltengrenze überall an derselben Stelle (gemeinsam.md „Satz“)
\newsavebox{\gbox}\newlength{\gLinks}
% v0.10 (Lehrer 13.10.): keine starre Halbierung – die linke Spalte ist so breit wie die Figur (+ 6 mm), die
% Antwort steht direkt daneben; ist die Figur breiter als 0,62 Zeile, steht die Antwort darunter
\newcommand{\gzwei}[2]{\par\smallskip\noindent\sbox{\gbox}{#1}%
  \setlength{\gLinks}{\dimexpr\wd\gbox+6mm\relax}%
  \ifdim\gLinks>0.62\linewidth
    \usebox{\gbox}\par\smallskip\noindent\begin{minipage}{\linewidth}#2\end{minipage}%
  \else
    \begin{minipage}[t]{\gLinks}\vspace{0pt}\usebox{\gbox}\end{minipage}%
    \begin{minipage}[t]{\dimexpr\linewidth-\gLinks\relax}\vspace{0pt}\raggedright #2\end{minipage}%
  \fi\par}
% v0.9 (12.10., pruefung.md § 5 „Marke mit B, Bedingungen rechts“): drittes Argument von pfaufg = Bedingungen
% der Prüfung (v0.10: \pfohnetr = TR durchgestrichen, „FOR“), klein grau im rechten Rand auf der Grundlinie der Nummer, Zeilen mit \\
\newcommand{\gbedrand}[2]{\ifblank{#2}{}{\rlap{\hspace*{\dimexpr#1+2mm\relax}%
  \parbox[t]{15mm}{\raggedright\scriptsize\leavevmode\color{mbgrau}\strut#2}}}}
% Teil eines Bündels (a, b …): Marke links im Rand, Bedingungen rechts im Rand, auf der Zeile des Teils
\newcommand{\gteilrand}[2]{\llap{\mbpfrandmarke{#1}\hspace*{\mbpfnr}}\gbedrand{\linewidth}{#2}}
% Teil eines Bündels: Buchstabe in eigener Spalte, Inhalt eingerückt (\gEin = Einzug, damit die Spaltengrenze von
% \gzwei an derselben Stelle bleibt)
\newenvironment{gbteil}[3]{\par\smallskip\noindent\gteilrand{#2}{#3}\makebox[\mbpfteil][l]{\textbf{#1}}%
  \begin{minipage}[t]{\dimexpr\linewidth-\mbpfteil\relax}\setlength{\gEin}{\mbpfteil}\raggedright\strut\ignorespaces}%
 {\end{minipage}\par}
% Grafik höchstens so breit wie die Spalte (Bündel nebeneinander, Koordinatensystem ohne Karo)
\newsavebox{\gpbox}
\newcommand{\gpassend}[1]{\sbox{\gpbox}{#1}\ifdim\wd\gpbox>\linewidth\typeout{GPASSEND: verkleinert \the\wd\gpbox\space auf \the\linewidth}\resizebox{\linewidth}{!}{\usebox{\gpbox}}\else\usebox{\gpbox}\fi}
% Skizzenfeld (v0.10, Lehrer 13.10.): Karo ab dem linken Rand über die Zeile, „Skizze“ klein grau oben links im Feld
% (wie \pfskizzenfeld der Vorlage); der rechte Rand bleibt den Bedingungen
\newcommand{\gskizze}[1]{\par\smallskip\noindent\begin{tikzpicture}
  \clip (0,0) rectangle ({\linewidth-1mm},{#1*10mm});
  \draw[black!16,line width=0.3pt,step=5mm] (0,0) grid ({\linewidth},{#1*10mm});
  \node[anchor=north west,inner sep=2pt,font=\scriptsize,text=mbgrau,fill=white] at (0,{#1*10mm}) {Skizze};
  \end{tikzpicture}\par}
\RenewDocumentEnvironment{pfaufg}{m m m}{%
  \begin{lrbox}{\mbaufgabebox}\begin{minipage}[t]{\linewidth}%
  \gSchrift\noindent\hspace*{\gEin}\mbpfrandmarke{#1}\gbedrand{\linewidth-\gEin}{#3}%
  \begin{minipage}[t]{\mbpfnr}\strut\textbf{#2}\end{minipage}%
  \begin{minipage}[t]{\dimexpr\linewidth-\gEin-\mbpfnr\relax}\raggedright\strut\ignorespaces}%
 {\end{minipage}%
  \end{minipage}\end{lrbox}%
  \setlength{\mbaufgabehoehe}{\dimexpr\ht\mbaufgabebox+\dp\mbaufgabebox\relax}%
  \ifdefempty{\gkopfnach}{\par\addvspace{5pt}\Needspace*{\mbaufgabehoehe}}{%
    % Stufenkopf, der mit dieser (ersten) Aufgabe zusammenbleibt: Platz für beide verlangen, dann den Kopf setzen
    % (12.10., Quadratische: ein Kopf stand allein unten, die Aufgabe mit Koordinatensystem auf der nächsten Seite)
    \par\Needspace*{\dimexpr\mbaufgabehoehe+40pt\relax}\gkopfnach\gdef\gkopfnach{}\par\addvspace{5pt}}%
  \noindent\usebox{\mbaufgabebox}\par\addvspace{5pt}}
% Stufenkopf vormerken: \gkopfvor{<Kopf>} – die nächste pfaufg setzt ihn, sobald beide zusammen Platz haben
\def\gkopfnach{}
\newcommand{\gkopfvor}[1]{\gdef\gkopfnach{#1}}
% Teilaufgabe mit Skizze links in gzwei: Buchstabe in eigener Spalte, darunter Text und Skizze
\newcommand{\gteil}[2]{\makebox[\mbpfteil][l]{\textbf{#1}}\parbox[t]{\dimexpr0.5\textwidth-\gEin-\mbpfnr-5mm-\mbpfteil\relax}{\raggedright #2}}
% Stufenkopf der Vorlage (mathblatt.sty 2026-10-12: Linie, Abstand, 13 pt, \pfjahre) ohne ihren Vorhalt von
% 14 Zeilen: der Setzer fordert vor dem Kopf 0,16 Seitenhöhe (Kopf + erste Aufgabe), mehr nicht
\setlength{\mbpfstufevorhalt}{0pt}
\makeatother
\newcommand{\karo}[1]{\par\smallskip\noindent\begin{tikzpicture}
  \clip (0,0) rectangle ({\linewidth-1mm},{#1*10mm});
  \draw[black!16,line width=0.3pt,step=5mm] (0,0) grid ({\linewidth},{#1*10mm});
  \end{tikzpicture}\par}
\newcommand{\kk}[1]{$\square$\,#1\hspace{0.9em}}
\newcommand{\linie}[1]{\,{\color{black!45}\rule[-1pt]{#1}{0.4pt}}\,}
\newcommand{\pkt}{\textperiodcentered{}}
% v0.14 (14.10.): Stufenkopf mit fünf Punkten (Option = Jahre für die Bestellkarte) oder „zuletzt 19“
\NewDocumentCommand{\gpunkte}{o m}{#2}
\newcommand{\gzuletzt}[1]{zuletzt #1}
% v0.7 (12.10.): Schülerfrage unter dem Stufenkopf; Ankreuzoption als Bild (Buchstabe und Kästchen darunter)
\newcommand{\gfrage}[1]{\par\nobreak\noindent\hspace*{\mbpfnr}{\small\itshape #1}\par\nobreak\addvspace{2pt}}
\newlength{\gbildbreite}\setlength{\gbildbreite}{0.32\linewidth}% Setzer setzt sie je Reihe (vier Bilder: 0,24)
\newcommand{\gbildoption}[3]{\begin{minipage}[t][#3][b]{\gbildbreite}\centering #2\par\smallskip\mbox{#1\,$\square$}\end{minipage}}
% Lineare (12.10.): Bild mit Buchstabe ohne Kästchen (Diagramm zuordnen); Aussage mit „wahr / falsch“ zum Ankreuzen
\newcommand{\gbildname}[3]{\begin{minipage}[t][#3][b]{\gbildbreite}\centering #2\par\smallskip\mbox{\textbf{#1}}\end{minipage}}
\newcommand{\gwahrfalsch}[2]{\par\noindent\hangindent2.2em\hangafter1\makebox[2.2em][l]{(#1)}#2\hfill\mbox{wahr $\square$\quad falsch $\square$}\par}
\makeatletter\setlength{\mbpfbe}{0mm}\makeatother
\begin{document}

\pfheftstil
\pflaufkopf[P10 \pkt{} mehr]{Volumen und Oberfläche}{PG90-6}\pfseitevon
\pfheftkopf{Volumen und Oberfläche}{mehr}
\pfstufenzeile{Volumen gesucht \pkt{} Oberfläche gesucht \pkt{} Körper vergleichen}
% 2026-FOR-B1f-m1 (Schwere I, rang 1)
\begin{pfaufg}{nach~P10\\’26~\pkt{}~1f}{1.}{}
Ein Würfel hat die Kantenlänge a = 4 cm. Gib sein Volumen an.
\karo{2}
\fusshilfe{\mbox{1:~64 cm³}}
\end{pfaufg}
% 2026-FOR-K2a-m1 (Schwere I, rang 2)
\begin{pfaufg}{nach~P10\\’26~\pkt{}~2a}{2.}{}
Ein Turm besteht unten aus einem Zylinder, sein Dach ist ein Kegel. Bekannt sind: h = 22,0 m (Höhe des Zylinders), r = 4,2 m (Radius von Zylinder und Kegel) und s = 7,9 m (Mantellinie des Kegels). Berechne das Volumen des zylinderförmigen Teils.
\par\smallskip\noindent\gpassend{\begin{tikzpicture}[baseline=(current bounding box.north),line width=0.7pt]\draw (-0.9,0) arc (180:360:0.9 and 0.315);\draw[dashed,line width=0.5pt] (-0.9,0) arc (180:0:0.9 and 0.315);\draw (-0.9,0) -- (-0.9,2.4);\draw (0.9,0) -- (0.9,2.4);\draw (0,2.4) ellipse (0.9 and 0.315);\draw (-0.9,2.4) -- (0,3.9) -- (0.9,2.4);\node[right,font=\small] at (0.95,1.2) {$h$};\fill (0,0) circle (0.8pt);\draw[dashed,line width=0.5pt] (0,0) -- (0.9,0) node[midway,fill=white,inner sep=1pt,font=\small] {$r$};\path (0.9,2.4) -- (0,3.9) node[midway,above right,inner sep=1.5pt,font=\small] {$s$};\end{tikzpicture}}\par
\karo{2}
\fusshilfe{\mbox{2:~1219,2 m³}}
\end{pfaufg}
% 2022-OS-K2b-m1 (Schwere I, rang 4)
\begin{pfaufg}{nach~P10\\’22~\pkt{}~2b}{3.}{}
Ein Becher hat die Form eines Zylinders ohne Deckel. Er ist h = 7,5 cm hoch und hat den Radius r = 3,4 cm. Zeige, dass 250 ml Flüssigkeit in den Becher passen (1 cm³ = 1 ml).
\par\smallskip\noindent\gpassend{\begin{tikzpicture}[baseline=(current bounding box.north),line width=0.7pt]\fill[black!12] (0,2) ellipse (0.8 and 0.28);\draw (0,2) ellipse (0.8 and 0.28);\draw (-0.8,0) -- (-0.8,2);\draw (0.8,0) -- (0.8,2);\draw (-0.8,0) arc (180:360:0.8 and 0.28);\draw[dashed,line width=0.5pt] (-0.8,0) arc (180:0:0.8 and 0.28);\end{tikzpicture}}\par
\karo{2}
\fusshilfe{\mbox{3:~272,4 cm³}}
\end{pfaufg}
% 2021-OS-K4a-m1 (Schwere I, rang 10)
\begin{pfaufg}{nach~P10\\’21~\pkt{}~4a}{4.}{}
Herr Gärtner stellt eine Regentonne in Form eines Zylinders auf. Sie ist h = 105 cm hoch und hat den Radius r = 27 cm. Der Hersteller gibt an, dass etwa 240 Liter Wasser in die Tonne passen. Bestätige diese Angabe mit einer Rechnung (1 l = 1 dm³).
\par\smallskip\noindent\gpassend{\begin{tikzpicture}[baseline=(current bounding box.north),line width=0.7pt]\draw (0,2.4) ellipse (0.8 and 0.28);\draw (-0.8,0) -- (-0.8,2.4);\draw (0.8,0) -- (0.8,2.4);\draw (-0.8,0) arc (180:360:0.8 and 0.28);\draw[dashed,line width=0.5pt] (-0.8,0) arc (180:0:0.8 and 0.28);\draw[line width=0.5pt] (-0.8,0.6) arc (180:360:0.8 and 0.28);\draw[line width=0.5pt] (-0.8,1.8) arc (180:360:0.8 and 0.28);\end{tikzpicture}}\par
\karo{2}
\fusshilfe{\mbox{4:~240,5 dm³}}
\end{pfaufg}
% 2018-OS-K6b-m1 (Schwere I, rang 11)
\begin{pfaufg}{nach~P10\\’18~\pkt{}~6b}{5.}{}
Ein runder Turm hat 5,8 m Durchmesser. Seine zylinderförmige Außenmauer ist 19,5 m hoch und soll renoviert werden. Berechne die Fläche der Außenmauer.
\karo{2}
\fusshilfe{\mbox{5:~355,3 m²}}
\end{pfaufg}
% 2014-OS-K5a-m1 (Schwere I, rang 12)
\begin{pfaufg}{nach~P10\\’14~\pkt{}~5a}{6.}{}
Für ein Café wird eine Kühlvitrine in Form eines Prismas gebaut. Die Grundfläche ist ein gleichschenkliges Trapez mit den parallelen Seiten 108 cm und 84 cm, den Schenkeln 37 cm und dem Flächeninhalt 3 360 cm². Die Vitrine ist 150 cm hoch. Berechne das Volumen der Vitrine.
\par\smallskip\noindent\gpassend{\begin{tikzpicture}[baseline=(current bounding box.north),x={(1cm,0cm)},y={(0.5cm,0.5cm)},z={(0cm,1cm)},line width=0.7pt]\draw (0,0,0) -- (3,0,0) -- (3,0,2.6) -- (0,0,2.6) -- cycle;\draw (3,0,0) -- (2.65,1.3,0) -- (2.65,1.3,2.6) -- (3,0,2.6);\draw (0,0,2.6) -- (0.35,1.3,2.6) -- (2.65,1.3,2.6);\draw[dashed,line width=0.5pt] (0,0,0) -- (0.35,1.3,0) -- (2.65,1.3,0);\draw[dashed,line width=0.5pt] (0.35,1.3,0) -- (0.35,1.3,2.6);\node[below,font=\small] at (1.5,0,0) {$108$};\node[above,font=\small] at (1.5,1.3,2.6) {$84$};\node[below right,font=\small] at (2.85,0.65,0) {$37$};\node[left,font=\small] at (0,0,1.3) {$150$};\end{tikzpicture}}\par
\karo{2}
\fusshilfe{\mbox{6:~504000 cm³}}
\end{pfaufg}
% 2021-OS-K4b-m1 (Schwere II, rang 13)
\begin{pfaufg}{nach~P10\\’21~\pkt{}~4b}{7.}{}
Herr Gärtner stellt eine Regentonne in Form eines Zylinders auf. Sie ist h = 105 cm hoch und hat den Radius r = 27 cm. Herr Gärtner meint: „Eine Tonne mit dreifachem Radius würde genau dreimal so viel Wasser fassen.“ Entscheide, ob das stimmt. Begründe deine Entscheidung.
\karo{2}
\fusshilfe{\mbox{7:~nein, Volumen verneunfacht sich: 240,5 l statt 2164,3 l}}
\end{pfaufg}
% 2016-OS-K3c-m1 (Schwere II, rang 14)
\begin{pfaufg}{nach~P10\\’16~\pkt{}~3c}{8.}{}
Eine Kugel zum Kugelstoßen hat einen Durchmesser von 11 cm. Berechne ihr Volumen.
\karo{2}
\fusshilfe{\mbox{8:~696,9 cm³}}
\end{pfaufg}
\par\vfill\noindent{\scriptsize\color{mbgrau}Originale: P10 2026 \pkt{} 1f, P10 2026 \pkt{} 2a, P10 2022 \pkt{} 2b, P10 2021 \pkt{} 4a, P10 2018 \pkt{} 6b, P10 2014 \pkt{} 5a, P10 2021 \pkt{} 4b, P10 2016 \pkt{} 3c\par}
\end{document}